\section{Lesson 8: Can We Demonstrate It?}

In Lesson 7 you interrogated pages of Aquinas, Scotus, and Hume unaided. Now the pressure reverses: attempt a full metaphysical demonstration from start to finish, judge honestly whether it succeeded, and state the weaker result when it did not. Constructing a proof makes every shortcut tempting---contested principles look obvious, definitions quietly widen, rivals shrink, and revealed truths borrow the voice of natural reason.

This final lesson adds no new question and no new rule. It holds bounding, denial classification, revealed-premise discipline, definition versus discovery, and rivals at full strength together under the pressure of building and judging. Our six headings remain a modern reading habit, never an outline taken from Aristotle or St.\ Thomas Aquinas; here we test real reconstructions directly against Aristotle's demonstrative standards (\emph{Posterior Analytics} I.2, 71b17--72a8; I.13, 78a22--79a16) and Aquinas's rules for proving causes from effects (\emph{Summa Theologiae} I, q.\,2, a.\,2).

\subsection{Worked Passage 1: Attempting a Full \emph{Demonstratio Quia} of a First Mover and First Efficient Cause}

Because Aquinas never gathers his demonstrative method and causal arguments into one paragraph, this passage explicitly reconstructs three verified parts of his corpus: \emph{ST} I, q.\,2, a.\,2, corp.\ and ad 2 (demonstration from effects); \emph{ST} I, q.\,2, a.\,3, corp.\ (the First and Second of his \emph{quinque viae}---``five ways,'' not ``Five Proofs''); and \emph{Summa contra Gentiles} I, c.\,13, nn.\,14, 15, and 33 (ordered movers, instrumental movers, and ordered efficient causes).

Begin with Aquinas's methodological charter in \emph{ST} I, q.\,2, a.\,2, ``Whether It Can Be Demonstrated That God Exists?'' (translated by the Fathers of the English Dominican Province, 1920):

\begin{quote}
Demonstration can be made in two ways: One is through the cause, and is called \emph{a priori}, and this is to argue from what is prior absolutely. The other is through the effect, and is called a demonstration \emph{a posteriori}; this is to argue from what is prior relatively only to us. When an effect is better known to us than its cause, from the effect we proceed to the knowledge of the cause. And from every effect the existence of its proper cause can be demonstrated, so long as its effects are better known to us; because since every effect depends upon its cause, if the effect exists, the cause must pre-exist. Hence the existence of God, in so far as it is not self-evident to us, can be demonstrated from those of His effects which are known to us.
\end{quote}

Check the Leonine Latin (Rome, 1888): where the 1920 English writes \emph{a priori} and \emph{a posteriori}, Aquinas writes \emph{propter quid} (\emph{per priora simpliciter}) and \emph{demonstratio quia} (\emph{per ea quae sunt priora quoad nos})---Aristotle's distinction (\emph{Post.\ An.}\ I.13, 78a22--79a16), not Kant's. In Reply to Objection 2, Aquinas answers the Aristotelian objection that demonstration requires the essence (\emph{quod quid est}) as middle term:

\begin{quote}
When the existence of a cause is demonstrated from an effect, this effect takes the place of the definition of the cause in proof of the cause's existence. This is especially the case in regard to God, because, in order to prove the existence of anything, it is necessary to accept as a middle term the meaning of the word, and not its essence, for the question of its essence follows on the question of its existence. Now the names given to God are derived from His effects; consequently, in demonstrating the existence of God from His effects, we may take for the middle term the meaning of the word ``God''.
\end{quote}

Aquinas substitutes \emph{quid significet nomen} (what the name signifies) for \emph{quod quid est} (the essence), since \emph{whether} a thing is precedes \emph{what} it is (\emph{Post.\ An.}\ II.8, 93a16--24; \emph{Super Boetium De Trinitate} q.\,6, a.\,4, ad 2). Applied to the First and Second Ways (\emph{ST} I, q.\,2, a.\,3, corp.) and backed by \emph{SCG} I, c.\,13, nn.\,14, 15, and 33, the reconstructed proof runs in five steps: (1) some sensible things undergo motion and stand in orders of efficient causes; (2) motion reduces potency to act, and nothing reduces itself from potency to act or causes itself efficiently, so whatever is moved or caused depends on another; (3) in an ordered series where subordinate members act only as moved or caused by a prior member (staff by hand), endless regress removes any first principal mover or first cause, eliminating intermediate causes and final effects, which is false (\emph{SCG} I, c.\,13, nn.\,14, 15, 33); (4) where a cause's essence is unknown, \emph{demonstratio quia} proves its existence from effects via \emph{quid significet nomen}; (5) therefore a first unmoved mover and first uncaused efficient cause exist---what everyone names ``God.''

\subsubsection*{Claim}

From sensible motion and ordered efficient causes, natural reason demonstrates \emph{quia}---taking \emph{quid significet nomen} as middle term---that a first unmoved mover and first uncaused efficient cause exist. \emph{Plain-language reading.} Observing change and ordered causal chains, and granting that nothing actualizes its own unrealized capacity or produces itself, we prove \emph{that} an unmoved, uncaused source heads concurrently dependent chains---not by grasping its inner nature, but by reasoning from effects under the functional meaning attached to ``God.'' We keep \emph{demonstratio quia}, \emph{quid significet nomen}, and \emph{quod quid est} because casual English erases Aquinas's substitution of nominal meaning for an unknown essence.

\subsubsection*{Meaning}

Five terms stay fixed: \emph{demonstratio quia} (proving \emph{that} from effects better known to us) versus \emph{propter quid} (showing \emph{why} from a cause better known in nature); \emph{quid significet nomen} (functional meaning from effects) versus \emph{quod quid est} (defining essence); \emph{motion} (\emph{motus}, reduction from real potency to act---alteration like wood heating, substantial change, or local change---not uniform inertial drift); \emph{ordered movers} and \emph{ordered efficient causes} (\emph{SCG} I, c.\,13, nn.\,14, 15, 33: \emph{per se} chains of concurrent derived power, distinct from \emph{per accidens} succession in time, \emph{ST} I, q.\,46, a.\,2, ad 7); and \emph{middle term} (here, the nominal account drawn from effects).

\subsubsection*{Warrant}

Keep \textbf{revelation} and bare authority apart from natural premises. The \emph{sed contra} of \emph{ST} I, q.\,2, a.\,2 cites Romans 1:20, and a.\,3 cites Exodus 3:14: both stand labelled as \textbf{revelation}, doing no philosophical work in the \emph{corpus}. In q.\,2, a.\,2, ad 1, Aquinas calls God's existence a preamble to the articles (\emph{praeambula fidei})---reachable by natural reason yet rightly held on faith by one who cannot follow a proof (\emph{SCG} I, c.\,3). Removing every biblical and Aristotelian name leaves the natural force unchanged; citing ``as Aquinas proved'' without the reasoning stands on \textbf{provenance only}---a diagnosis, never a warrant.

Premise by premise: step 1 rests on \textbf{experience} (\emph{certum est enim, et sensu constat}). Step 2 blends \textbf{definition} (change as passage from potency to act) with substantive \textbf{metaphysical principle} (real potency requiring concurrent actualization by something in act, joined to non-contradiction and the \textbf{metaphysical principle} barring self-causation). Step 3 enters \emph{ST} I, q.\,2, a.\,3 as a \textbf{prior result} argued in \emph{SCG} I, c.\,13, nn.\,14, 15, and 33 (ordered movers, n.\,14; instrumental movers, n.\,15; ordered efficient causes, n.\,33)---standing on \textbf{provenance only} inside the \emph{Summa} until checked, where it rests on a \textbf{metaphysical principle} about subordinate agency illustrated from \textbf{experience} by hand and staff. Step 4 pairs \textbf{definition} of the demonstrative modes with the \textbf{metaphysical principle} that every effect depends on its proper cause. Step 5 follows by \textbf{logical consequence}, attaching ``God'' by nominal \textbf{definition}.

\subsubsection*{Burden}

Ask the hostile-reader question: \emph{could an intelligent opponent accept every definition here and still deny the conclusion?} Yes---isolating two linked burden-bearing premises: first, \textbf{that potency is a real principle in changing things requiring concurrent actualization by an external agent in act, not merely a modal summary of physical laws}; second, \textbf{that actual natural causal orders are genuinely ordered \emph{per se}---concurrent hierarchies of derived acting beyond local hand-and-stick mechanisms---rather than \emph{per accidens} sequences of intrinsically powered agents in time.} Deny either bridge and the proof halts.

\subsubsection*{Consequence}

Before judging the proof, expose a structural ambiguity English buries inside ``every member depends on another.'' Here, and nowhere else in Lesson 8, two displayed lines earn their place because prose easily merges two genuinely different formal readings of a causal series. Fix a present moment \linline{t}, let quantifiers range over subordinate members of a series, let \linline{\logicop{Sust}(x, y, t)} read \emph{at \linline{t}, \linline{x}'s present acting depends concurrently on \linline{y}'s present acting}, and let \linline{\logicop{Prec}(x, y)} read \emph{\linline{x} was brought about by an earlier \linline{y}}. Compare:
\ldisplay{(1) \lquad \lall x\lthin \lex y\lthin \logicop{Sust}(x, y, t)}
\ldisplay{(2) \lquad \lall x\lthin \lex y\lthin \logicop{Prec}(x, y)}
This display exposes the exact inferential feature separating the two structures: line (1) captures the concurrent \emph{per se} series required by step 3, where each subordinate member borrows its present acting simultaneously at \linline{t} (so removing an underived source removes the present effect), whereas line (2) captures bare historical \emph{per accidens} succession in time, permitting endless regress without removing present activity. The symbols deliberately set aside \emph{whether} cosmic causes actually sustain one another concurrently (left in prose under Burden), strictly to keep concurrent \emph{per se} dependence visibly separate from temporal succession.

Now test the reconstruction systematically against Aristotle's six conditions on demonstrative premises (\emph{Posterior Analytics} I.2, 71b17--72a8) and the \emph{quia} / \emph{propter quid} distinction (\emph{Post.\ An.}\ I.13, 78a22--79a16). Strict demonstration (\emph{apodeixis}) demands premises that are \emph{true, primary, immediate, better known than the conclusion, prior to it, and causes of it}---``better known'' meaning better known \emph{in nature} (71b29--72a5):

\begin{itemize}
\item \emph{True:} conditional on the two contested metaphysical principles named under Burden.
\item \emph{Primary and immediate:} sensory awareness of change is immediate, yet neither real ontological potency nor cosmic \emph{per se} ordering is an immediate axiom known without argument.
\item \emph{Better known in nature:} explicitly \emph{not} met; Aquinas grants (\emph{ST} I, q.\,2, a.\,2, corp.) that we argue from effects better known \emph{to us}.
\item \emph{Prior to and causes of the conclusion:} explicitly \emph{not} met in reality; sensible motion is an effect, and the middle term is \emph{quid significet nomen} rather than \emph{quod quid est}.
\end{itemize}

State the verdict directly, without hedging and without apology: \textbf{at best, this reconstruction yields a conditional \emph{demonstratio quia}---conditional on contested metaphysical principles (potency as a real principle of change; real \emph{per se} ordered dependence in nature), using what the name signifies rather than the essence as its middle term.} It is strictly \emph{quia}, never \emph{propter quid}.

Whether even this \emph{quia} reconstruction achieves strict Aristotelian demonstrative force remains a genuine scholarly dispute. Critics such as Anthony Kenny (\emph{The Five Ways}, 1969) judge that the Ways fail as strict Aristotelian demonstrations, citing obsolete physical assumptions, contested causal rules, and gaps in identifying the terminus with God. By contrast, historians and defenders such as John F.\ Wippel (\emph{The Metaphysical Thought of Thomas Aquinas}, 2000) and Edward Feser (\emph{Five Proofs of the Existence of God}, 2017) argue that, read within Aquinas's fuller metaphysics of \emph{esse} and concurrent \emph{per se} ordering, the Ways sustain genuine \emph{quia} demonstrative force. Even among Aristotelians, because being is not a genus (\emph{Metaphysics} III.3, 998b22) and God's essence remains unknown to us, whether a \emph{quia} proof anchored in a nominal definition yields strict \emph{episteme} in the sense of \emph{Posterior Analytics} I.2 remains open. We do not settle that dispute by fiat.

A reader who grants every premise \emph{has} a conditional \emph{quia} inference establishing \emph{that} each real \emph{per se} ordered series ends in an unmoved or uncaused first source answering to the functional name ``God.''

\emph{Not yet established.} Granting every premise still does \emph{not} yield: (1) knowledge of God's essence (\emph{quod quid est}, explicitly disclaimed at \emph{ST} I, q.\,2, a.\,2, ad 2--3); (2) a \emph{propter quid} explanation of \emph{why} God or creatures exist; (3) uniqueness across distinct series; or (4) the divine attributes---simplicity, goodness, infinity, immutability, eternity, and unity, argued separately in \emph{ST} I, qq.\,3--11 (and \emph{SCG} I, cc.\,14--28), plus divine knowledge, argued separately in \emph{ST} I, q.\,14.

\subsubsection*{Denial}

Rejecting either burden-bearing premise is \textbf{rival framework} (paired with a live contest over \textbf{explanatory cost}), and specifically \emph{not} \textbf{contradiction} and \emph{not} \textbf{conflict with experience}. A denier affirms every law of logic and every sensible fact---fire heating wood, bodies colliding, organisms growing---without contradicting himself; calling denial ``incoherent'' overreaches. Simply declining the act/potency vocabulary is \textbf{non-assent}.

State the strongest rivals fairly, as their ablest defenders would accept them. First, against real potency, a developed microphysical or dispositional framework holds that ``wood is potentially hot'' summarizes modal regularities grounded in molecular structure and physical laws (including inertial motion persisting without a concurrent mover), not a distinct real principle needing continuous external actualization. Second, against cosmic \emph{per se} ordering, a physicalist or Humean framework---seen in Lesson 7 through Cleanthes's twenty particles in David Hume's \emph{Dialogues Concerning Natural Religion}, Part IX (1779)---maintains that nature exhibits temporal \emph{per accidens} succession (\linline{\lall x\lthin \lex y\lthin \logicop{Prec}(x, y)}) among intrinsically powered systems governed by conservation laws rather than a vertical ladder of borrowed agency (\linline{\lall x\lthin \lex y\lthin \logicop{Sust}(x, y, t)}). Sticks in hands are local tools, not proof that fields derive ongoing acting from a transcendent First Cause; even admitting vertical grounding layers, ruling out an endless downward stack requires independent proof.

In response, the Thomist presses a serious \textbf{explanatory cost}: without real potency and concurrent vertical dependence on an underived act of being, the rival leaves unexplained why contingent physical laws obtain, why finite things possess determinate capacities, and why contingent reality continues right now. That cost is weighty philosophical pressure---yet still an explanatory cost, not a contradiction; the rival may answer internally or accept fundamental laws as brute.

\subsection{Worked Passage 2: An Attempted Demonstration That Fails---Proving the Universe Began in Time}

Sometimes a serious historical candidate demonstration does not even yield a conditional proof: it fails outright. Honesty requires stating that verdict directly, refusing a disguised definitional repair, and stating cleanly what weaker truth survives.

Consider the classic medieval attempt to demonstrate by natural reason that the universe had a temporal beginning, recorded by Aquinas as Objections 6 and 7 of \emph{ST} I, q.\,46, a.\,2, ``Whether It Is an Article of Faith That the World Began?'' (1920 English Dominican translation):

\begin{quote}
Obj.\ 6: Further, if the world always was, the consequence is that infinite days preceded this present day. But it is impossible to pass through an infinite medium. Therefore we should never have arrived at this present day; which is manifestly false.

Obj.\ 7: Further, if the world was eternal, generation also was eternal. Therefore one man was begotten of another in an infinite series. But the father is the efficient cause of the son (Phys.\ ii, text 5). Therefore in efficient causes there could be an infinite series, which is disproved (Metaph.\ ii, text 5).
\end{quote}

Combined, Objections 6 and 7 attempt a \emph{reductio}: an eternal cosmos entails crossing infinitely many past days and an endless chain of efficient causes in generation; both are impossible, so reason proves the world began.

\subsubsection*{Claim}

Natural reason demonstrates that the universe had a temporal beginning, because an eternal past would require both traversing infinite bygone days to arrive at today and an infinite regress of efficient causes in human generation.

\emph{Plain-language reading.} If the universe had existed forever, an endless sequence of past days would have had to finish before today could arrive, and an endless chain of fathers would have had to beget sons; since neither is possible, philosophy proves the cosmos began a finite time ago.

\subsubsection*{Meaning}

Four terms must stay fixed: \emph{passing through an infinite medium} (crossing between two determinate endpoints versus looking backward along an unbounded series); \emph{infinite series in efficient causes} (\emph{per se} derived causing versus \emph{per accidens} temporal succession); \emph{eternal} applied to the cosmos (endless successive time, distinct from God's simultaneous timeless eternity, \emph{ST} I, q.\,46, a.\,2, ad 5); and \emph{creation} (\emph{creatio ex nihilo}, total ontological dependence---``not made from anything''---distinct from ``made after non-existence in time,'' ad 1--2).

\subsubsection*{Warrant}

Keep \textbf{revelation} strictly separate from the candidate proof. In \emph{ST} I, q.\,46, a.\,2, corp., Aquinas holds the conclusion by faith alone (\emph{mundum non semper fuisse, sola fide tenetur, et demonstrative probari non potest}), resting on Genesis 1:1 and the Creed; Scripture lends zero warrant to Objections 6 and 7, and their failure weakens nothing held by faith.

Inside the attempt: today's arrival and human generation rest on \textbf{experience}; inferring infinite past days is a \textbf{logical consequence} of unbounded past duration, whereas inferring infinite human generations slips in an unargued \textbf{hypothesis} that humanity existed eternally with the cosmos (ad 8). That an infinite medium cannot be traversed enters as a \textbf{metaphysical principle}. Objection 7's citations to \emph{Physics} II and \emph{Metaphysics} II supply \textbf{provenance only}; philosophically, the no-regress rule is a \textbf{prior result} valid strictly for \emph{per se} causes. The conclusion claims to follow by \textbf{logical consequence}.

\subsubsection*{Burden}

Two burden-bearing premises carry the attempt, both deniable while granting every definition: from Objection 6, \textbf{that an unbounded eternal past with no first day constitutes a traversal across an infinite medium between two fixed extremes}; and from Objection 7, \textbf{that successive generations form a \emph{per se} ordered series of efficient causes barred by the no-regress rule} (plus the hidden assumption that humanity is co-eternal with the world).

\subsubsection*{Consequence}

State the verdict in one crisp, direct sentence, without apology: \textbf{This attempt fails as a strict demonstration.}

Both middle terms collapse under Aquinas's replies (\emph{ST} I, q.\,46, a.\,2, ad 6 and ad 7). Against Objection 6, Aquinas answers verbatim (1920 English Dominican translation):

\begin{quote}
Passage is always understood as being from term to term. Whatever bygone day we choose, from it to the present day there is a finite number of days which can be passed through. The objection is founded on the idea that, given two extremes, there is an infinite number of mean terms.
\end{quote}

Objection 6's middle term silently assumes a first day infinitely far back; without a first day, any chosen bygone day stands only finitely many days before today. Against Objection 7, Reply 7 (Lesson 7) exposes an equivocation in ``infinite series of efficient causes'': \emph{SCG} I, c.\,13, n.\,33 bars endless \emph{per se} stacks (stone, stick, hand), not \emph{per accidens} lines on one causal tier (broken hammers, or fathers begetting \emph{inquantum homo} rather than \emph{inquantum filius}). Deeper still, q.\,46, a.\,2, corp.\ shows why no \emph{propter quid} proof exists: creaturely essences abstract from ``here'' and ``now,'' and God creates freely rather than by necessity.

Now exhibit and \textbf{explicitly refuse a tempting definitional rescue}: redefining ``created universe'' to mean \emph{a universe that began to exist after a state of non-existence in time}. Paired with \emph{ST} I, q.\,44, a.\,1 (the universe is created by God), a temporal beginning would follow immediately. We refuse this repair outright because it collapses a substantive cosmological question into vocabulary unpacking: any opponent defending an eternal yet continuously created cosmos grants the stipulated definition and simply declines to apply that temporal label to the universe, holding that ontological creation (\emph{ex nihilo}, not from pre-existing matter, ad 2) entails total dependence without requiring a first moment in time.

What weaker result genuinely survives? First, the inquiry exposes a real \textbf{dialectical and conceptual pressure}: even though bounded traversal fails against a beginless past, the live tension surrounding completed actual infinites in past time remains a serious philosophical problem (probed further in Obj.\ 8 on immortal souls). Second, it separates \textbf{ontological creation} cleanly from \textbf{temporal beginning}. Third, it leaves the world's temporal beginning honestly labelled as \textbf{revelation} (\emph{sola fide tenetur}, \emph{ST} I, q.\,46, a.\,2, corp.)---neither laundered into natural demonstration nor weakened as an article of faith when a non-cogent proof fails. Aquinas states the discipline unmistakably:

\begin{quote}
Hence that the world began to exist is an object of faith, but not of demonstration or science. And it is useful to consider this, lest anyone, presuming to demonstrate what is of faith, should bring forward reasons that are not cogent, so as to give occasion to unbelievers to laugh, thinking that on such grounds we believe things that are of faith.
\end{quote}

\emph{Not yet established.} Natural reason has demonstrated neither that the universe began in time nor that an eternal created cosmos is impossible---nor, symmetrically (\emph{ST} I, q.\,46, a.\,1, corp.), that the universe must be eternal, since God wills creatures freely.

\subsubsection*{Denial}

Rejecting both burden-bearing premises incurs neither \textbf{contradiction} nor \textbf{conflict with experience}; it stands on a \textbf{rival framework} whose ablest expositor here is St.\ Thomas Aquinas himself (\emph{ST} I, q.\,46, a.\,2, ad 1, ad 2, ad 6, ad 7).

Stated at full strength, that coherent modal framework is an \textbf{eternal yet continuously created cosmos}. Using the Platonists' image reported by St.\ Augustine (\emph{De Civitate Dei} X, 31, cited in ad 1), a foot eternally in dust would produce an eternal footprint depending instantaneously and wholly on the foot. Because creation is instantaneous ontological giving rather than successive motion from prior matter, an eternal cosmos could depend wholly on God without God preceding it in clock-time: every past day lies finitely far from today (ad 6), generations run \emph{per accidens} (ad 7), and creation \emph{ex nihilo} means ``not from anything'' rather than ``after nothing'' in time (ad 2).

\subsection{Where This Leaves You}

The scaffolding can now retire from the table. Keep no chapter titles, no rules to recite, and no proprietary apparatus. Beside any open page of Aristotle, Augustine, Aquinas, Scotus, or a modern metaphysician, carry only a pencil, blank paper, the One-Page Analysis Sheet (Appendix B) if you want a desk reminder, and six ordinary questions: What is claimed in plain language with every qualification intact? What do key terms mean here, and do any shift? Why accept each premise---definition, logical consequence, experience, prior result, metaphysical principle, hypothesis, revelation, or provenance only? Which sentence carries the substantive burden beyond vocabulary and deferred prior results? What actually follows, and what is not yet established? What does denial cost---contradiction, conflict with experience, explanatory cost, a rival framework at full strength, or non-assent?

Apply the discipline of bounding conclusions one last time, to metaphysics itself. \textbf{Neither of two sweeping verdicts has been established: metaphysical demonstration has been shown neither to be generally available nor to be generally impossible.} Both sweeping verdicts exceed the evidence: claiming strict demonstrations are routinely available ignores Aristotle's austere conditions, while claiming demonstration is impossible in principle smuggles in an unargued global scepticism. Leaving both sweeping verdicts unasserted is the final habit: test every argument on its own page, one sentence at a time.

\subsection{Exercises}

\subsubsection*{Exercise 8.1}
Here is an unannotated and previously unexamined passage from St.\ Thomas Aquinas, \emph{ST} I, q.\,44, a.\,2, corp., ``Whether Primary Matter Is Created by God?'' (1920 English Dominican translation):

\begin{quote}
The ancient philosophers gradually, and as it were step by step, advanced to the knowledge of truth. At first being of grosser mind, they failed to realize that any beings existed except sensible bodies. And those among them who admitted movement, did not consider it except as regards certain accidents, for instance, in relation to rarefaction and condensation, by union and separation. And supposing as they did that corporeal substance itself was uncreated, they assigned certain causes for these accidental changes, as for instance, affinity, discord, intellect, or something of that kind. An advance was made when they understood that there was a distinction between the substantial form and matter, which latter they imagined to be uncreated, and when they perceived transmutation to take place in bodies in regard to essential forms. Such transmutations they attributed to certain universal causes, such as the oblique circle [\emph{The zodiac, according to Aristotle (De Gener.\ ii)}], or ideas, according to Plato. But we must take into consideration that matter is contracted by its form to a determinate species, as a substance, belonging to a certain species, is contracted by a supervening accident to a determinate mode of being; for instance, man by whiteness. Each of these opinions, therefore, considered ``being'' under some particular aspect, either as ``this'' or as ``such''; and so they assigned particular efficient causes to things. Then others there were who arose to the consideration of ``being,'' as being, and who assigned a cause to things, not as ``these,'' or as ``such,'' but as ``beings.''

Therefore whatever is the cause of things considered as beings, must be the cause of things, not only according as they are ``such'' by accidental forms, nor according as they are ``these'' by substantial forms, but also according to all that belongs to their being at all in any way. And thus it is necessary to say that also primary matter is created by the universal cause of things.
\end{quote}

Reconstruct and judge this argument as a candidate demonstration unaided under all six questions (middle term, \emph{quia}/\emph{propter quid} standing, warrant of every step including the historical sketch, burden, bounded conclusion with what is not yet established, and denial).

\subsubsection*{Exercise 8.2}
Read this unannotated paragraph:

\begin{quote}
Every changing physical thing around us is contingent and depends on prior physical conditions for its present state. Nothing contingent explains its own existence, and a chain of contingent physical conditions cannot extend backward forever without leaving the entire cosmos unexplained. Therefore, by strict \emph{propter quid} demonstration, there necessarily exists a single eternal, simple, all-loving personal Creator who freely brought the universe into being from nothing at the beginning of time. Anyone who denies this rejects reason itself and falls into self-contradiction.
\end{quote}

Diagnose where this attempted demonstration fails, and \textbf{weaken the argument into the strongest genuine claim its actual premises support}, naming every unbridged gap.

\subsubsection*{Exercise 8.3}
Construct a concise argument in ordinary prose, without quoting Aquinas, for the claim \textbf{that a wholly material body cannot be the unmoved first mover of \emph{ST} I, q.\,2, a.\,3.} Audit your own construction unaided under all six questions, judge honestly whether it reaches strict demonstration or stops at a conditional or dialectical result, run the hostile-reader test on your burden-bearing premise, and exhibit one tempting definitional cheat that your construction refuses.

\subsection{Solutions}

\subsubsection*{Solution to Exercise 8.1}
\emph{Claim.} Primary matter is created by the universal cause of being, because whatever causes things insofar as they are \emph{beings}---not merely ``such'' (accident) or ``this'' (substantial form)---must cause everything pertaining to their being in any way. \emph{Plain-language reading.} Earlier thinkers explained only changes in a body's traits or kind, leaving raw matter assumed; once you ask what causes things simply to exist at all, that universal cause must produce everything entering their existence, including their underlying material principle.

\emph{Meaning.} Fix \emph{primary matter} (bare hylomorphic potentiality, not formed stuff like clay), \emph{being as ``such'' / ``this'' / ``being''} (accidental, substantial, existential tiers), and \emph{universal cause of things} (cause of things simply as beings). \emph{Warrant.} The historical sketch supplies \textbf{provenance only} as proof of truth---framing the conceptual ascent rather than proving creation. In paragraph two, matter's contraction by form rests on \textbf{definition} and hylomorphic \textbf{prior result}; the existence of a universal cause of beings enters as a \textbf{prior result} from \emph{ST} I, q.\,44, a.\,1 (\textbf{provenance only} here until audited alongside qq.\,3 and 11); that a cause of things \emph{as beings} produces every constituent of their being is a \textbf{metaphysical principle} nearing \textbf{logical consequence}; primary matter's creation follows by \textbf{logical consequence}.

\emph{Burden.} Beyond q.\,44, a.\,1, two burdens sit here: \textbf{that bodies genuinely contain a real principle of pure potentiality (primary matter), and that the universal cause of being produces passive co-principles themselves rather than ordering an independent passive substrate.} \emph{Consequence.} With \emph{belonging to the being of things caused as beings} as middle term, the inference runs downward from the universal cause (q.\,44, a.\,1) to one constituent of its effect---giving it the direction of a \textbf{conditional \emph{propter quid} corollary} relative to q.\,44, a.\,1, though overall it inherits the conditional \emph{quia} standing of q.\,2, a.\,3 and q.\,44, a.\,1 (both descriptions are defensible; neither yields an unconditional self-standing proof here). \emph{Not yet established:} (1) that primary matter exists apart from form (denied in Reply 3); (2) a temporal beginning (q.\,46); (3) free rather than necessary creation; and (4) the imported arguments of q.\,44, a.\,1 and q.\,3, a.\,4. \emph{Denial.} \textbf{Rival framework} (atomism or field physicalism rejecting prime matter, or craft-dualism accepting an \textbf{explanatory cost} over why two ultimate principles fit together), not \textbf{contradiction}.

\subsubsection*{Solution to Exercise 8.2}
\emph{Diagnosis.} The attempt fails at four points. First, ``depends on prior physical conditions'' describes a \emph{per accidens} temporal sequence yet invokes a no-regress ban without proving concurrent \emph{per se} dependence or answering Hume's Cleanthes on the member-to-whole step. Second, calling an effect-to-cause inference a \emph{strict propter quid demonstration} reverses causal priority. Third, one sentence jumps across unargued gaps: uniqueness (``single''), entitative attributes (``eternal, simple,'' \emph{ST} I, qq.\,3--11), personal attributes (``all-loving personal,'' qq.\,14--20), and ``freely at the beginning of time''---laundering \textbf{revelation} (\emph{ST} I, q.\,46, a.\,2) into natural philosophy. Fourth, calling denial ``self-contradiction'' misclassifies a dispute over \textbf{rival framework} and \textbf{explanatory cost}.

\emph{Honest weakening.} The strongest claim supported is conditional and bounded: \textbf{If every contingent physical state depends for its existence on something beyond itself, and if an endless chain of contingent conditions is impossible without an external ground, then at least one non-contingent ground of the physical series exists---reached at best by conditional \emph{demonstratio quia} or explanatory \emph{hypothesis}, not \emph{propter quid} proof.} (A stricter reader might weaken it further to a dialectical pressure on brute contingency; both respect the premises.) \emph{Unbridged gaps:} (1) why a temporal contingent chain cannot be brute; (2) why the non-contingent ground cannot be a necessary physical law or field; (3) uniqueness; (4) simplicity and eternity; (5) intellect, freedom, love, and personhood; and (6) a temporal beginning (remaining \textbf{revelation}).

\subsubsection*{Solution to Exercise 8.3}
\emph{Model construction.} Assume conditionally from Worked Passage 1 that an unmoved first mover exists---a source of change not itself reduced from potency to act in the causal order it heads. Every wholly material body is spatially extended, composed of parts, open to division or alteration, and in experience moves other bodies only through reciprocal physical interaction. Hence every wholly material body carries real passive potentiality and dependence on external conditions. Since an unmoved first mover cannot stand in passive potency regarding the motion it originates, no wholly material body can be the unmoved first mover.

\emph{Self-audit.} \emph{Claim and plain-language reading.} No wholly material body can be the unmoved first mover of \emph{ST} I, q.\,2, a.\,3, because corporeal extension and interaction entail passive potentiality incompatible with an unmoved first source. \emph{Meaning.} Fix \emph{wholly material body} (an extended physical entity with parts), \emph{unmoved first mover}, and \emph{passive potentiality} (real capacity to be divided or altered). \emph{Warrant.} The first mover enters as a conditional \textbf{prior result} (inheriting Worked Passage 1's burdens); divisibility and reciprocal pushing join \textbf{experience} with a \textbf{metaphysical principle} linking extension to potency (compare \emph{ST} I, q.\,3, a.\,1, corp.); extending ``unmoved in this order'' to ``wholly free of passive potency'' adds a further \textbf{metaphysical principle}; the exclusion follows by \textbf{logical consequence}. \emph{Burden.} Beyond q.\,2, a.\,3, the burden is \textbf{that the unmoved first mover must be wholly devoid of passive potentiality (not merely unmoved in one respect), and that every wholly material reality necessarily involves real passive potentiality.}

\emph{Consequence and judgment.} This construction stops short of unconditional \emph{propter quid} demonstration; at best it yields a \textbf{conditional \emph{quia} (negative corollary) argument}, saying what the first mover is \emph{not} (\emph{SCG} I, c.\,14) rather than grasping its essence. \emph{Not yet established:} uniqueness; intellect, consciousness, or personhood (versus an impersonal immaterial principle); goodness; or creation from nothing. \emph{Denial.} \textbf{Rival framework}, not \textbf{contradiction}: a serious physicalist argues that a fundamental physical reality (such as an indivisible quantum field or invariant spacetime law) need not behave like a classically divisible, pushed solid yet remains wholly physical, or that a body unmoved in one respect need not be pure act in every respect. \emph{Hostile-reader test and refused cheat.} An opponent can grant both fixed definitions (``extended physical entity'' and ``source of change not itself moved in that order'') and still deny the conclusion, proving the burden is substantive. By contrast, a \textbf{tempting definitional cheat} defines ``material body'' as ``anything incapable of moving unless moved by another'' (or ``unmoved mover'' as ``a purely spiritual substance''): any physicalist accepts that stipulation and denies that fundamental physical reality fits the rigged label, leaving reality untouched.
